

% Begin insert: directly after \begin{document}

\pagestyle{empty}
\newcommand\sphinxbackoftitlepage{


Welcome to the documentation of \sphinxstylestrong{XlCalcNet}, an addin for MS E\sphinxstylestrong{x}ce\sphinxstylestrong{l} and LibreOffice \sphinxstylestrong{Calc}, based on Python\sphinxstylestrong{Net}. This edition of the documentation describes release 1.0.


\sphinxstylestrong{XlCalcNet} is a Python library for mathematical functions in multiprecision, based on mpmath (see \sphinxurl{http://mpmath.org}). It is free software released under the Mozilla Public License 2.0 (see {\hyperref[\detokenize{900_Back_Matter/10_License:rst-bsd-3-clause-license}]{\sphinxcrossref{\DUrole{std,std-ref}{License}}}}).


The git repository is \sphinxurl{https://github.com/duhadler}.
There is an issue tracker \sphinxurl{https://github.com/duhadler} for bug reports and feature requests.

This documentation is available in PDF format at \sphinxurl{https://duhadler.files.wordpress.com/2017/06/mpformula.pdf},
and in HTML format at \sphinxurl{http://mpformula.github.io/C/},


\vspace{18.0cm}
\noindent
Copyright \copyright \hspace{0.5mm}  2026 Dietrich Hadler

}\sphinxmaketitle
\pagestyle{plain}

% Tex for Preface: only one page!
\pagenumbering{roman}%
\chapter*{Preface}
\addcontentsline{toc}{chapter}{Preface}

\sphinxstylestrong{XlCalcNet} is a Python  library focussing on the calculation of mathematical functions in arbitrary
precision. It uses the mpmath library (see \sphinxurl{http://mpmath.org})  as its computational backend, and
is expected to work on all Python configurations on Windows, macOS and Linux which can run mpmath.

\sphinxstylestrong{XlCalcNet} is intended to be used together with existing software for numerical computing, in particular
the Python libraries NumPy, SciPy and Matplotlib, but also recent versions of RStudio and R, MATLAB, 
Microsoft Excel and LibreOffice Calc. It is also possible to interact with .NET languages
(like C\# or Visual Basic) to speed up calculations.

Examples for how to do this are given in the introduction, and comparisons of speed and accuracy of
existing numerical software are given in the documentation of many functions.


\newpage

%This page intentionally left blank

%\newpage



\sphinxtableofcontents
\pagestyle{normal}
\phantomsection\label{\detokenize{index::doc}}


% End insert: directly before \part{Getting started}



